% !PRINTED 408
\weberpage{408}
%<<CONT
rationaler Functionen von $x_0, x_1, \ldots, x_{n-1}$, die durch die
Permutationen der alternirenden Gruppe aus einer von ihnen hervorgehen und
Wurzeln der Gleichung

\weq{1}{F\,(u, z) = 0}

sind, worin $z$ eine rationale Function der $x$ ist, die durch die
Permutationen der alternirenden Gruppe ungeändert bleibt.

Wir beweisen folgenden zweiten Hülfssatz:

\gesperrt{2. Wenn die rationale Function $\Psi\,(u, u_1, \ldots, u_{\nu-1})$
identisch verschwindet, wenn für die
\[
x_0, x_1, \ldots, x_{n-1}
\]
gewisse rationale Functionen einer Variablen $t$}

\weq{2}{x_h = \varphi_h\,(t)}

\gesperrt{gesetzt werden, und wenn durch diese Substitution $z$ nicht von
$t$ unabhängig wird, so verschwindet $\Psi$ identisch auch als Function der
unabhängigen Variablen $x_0, x_1, \ldots, x_{\nu-1}$.}
% [sic] The source prints x_{\nu-1} here, although the independent variables were introduced through x_{n-1}.

Die Galois'sche Gruppe der Gleichung (1) in dem Körper der rationalen
Functionen von $z$ wird aus gewissen Permutationen der Wurzeln
$u, u_1, \ldots, u_{\nu-1}$ bestehen. Führen wir diese Permutationen in
der Function $\Psi$ aus, und bilden das Product

\weq{3}{\Pi\,\Psi\,(u, u_1, \ldots, u_{\nu-1}) = f\,(z)}

aller so erhaltenen Functionen, so ergiebt sich eine rationale Function
$f\,(z)$ von $z$. Wenn nun einer der Factoren des Productes (3) nach der
Substitution (2) identisch verschwindet, und $z$ ist nicht von $t$
unabhängig, so muss $f\,(z)$ identisch gleich Null sein, und folglich muss
einer der Factoren des Productes (3) auch als Function der $x$ identisch
verschwinden. Wenn aber einer dieser Factoren identisch gleich Null ist,
so verschwinden auch alle anderen Factoren, weil man in jeder rationalen
Gleichung zwischen den $u$ alle Permutationen der Galois'schen Gruppe
ausführen kann. Damit ist der Satz 2. bewiesen.

Dieser Satz giebt nun mit dem im vorigen Paragraphen bewiesenen Satze 1.
zusammengenommen das folgende Resultat:

\gesperrt{3. Sind $u, u_1, \ldots, u_{\nu-1}$ die Wurzeln einer von einem
Parameter $z$ abhängigen Gleichung (1), und sind $u, u_1, \ldots,
u_{\nu-1}, z$ rationale Functionen der}
%CONT>>
